\chapter{球面、双曲空間、そして\(k\)の意味}

少し数学をやろう、簡単な数学だ。任意の次元の球面は特別な幾何学である。すなわち一様であり、ここで一様とは曲率が一定であることを意味する。まず正の曲率、次に一様な負の曲率、そしてその後でアインシュタイン方程式との関係が来る。

球状宇宙とは、時空ではなく\emph{空間}が球状であることを意味する。もう一つの可能性、一様な負の曲率も存在する。これらの場合を並べて見ることで、一般相対性理論の役割が明らかになる。

\section{球面と一様幾何学}

一次元球面は単なる円である。それを記述する最も簡単な方法は、二次元の補助空間に埋め込んで
\begin{equation}
x^2+y^2=R^2.
\end{equation}
と書くことである。単位一次元球面は
\begin{equation}
x^2+y^2=1.
\end{equation}
である。

二次元球面も同じように記述されるが、今度は三つの補助座標を用いる：
\begin{equation}
x^2+y^2+z^2=R^2,
\end{equation}
単位の場合には
\begin{equation}
x^2+y^2+z^2=1.
\end{equation}
となる。

これらの座標は単なる補助にすぎない。関心のある幾何学は周囲の\(x,y,z\)空間ではない。曲面そのものである。一次元球面上に住む生物は円の上に住んでおり、内側にも外側にもアクセスできない。同じ考えが二次元球面にも当てはまる。球面の内部も球面の外部も、補助的な描像に属するだけである。

同じ構成で三次元球面が得られる。もう一つ座標、たとえば\(w\)を導入して
\begin{equation}
x^2+y^2+z^2+w^2=R^2,
\end{equation}
と書く。単位の場合には
\begin{equation}
x^2+y^2+z^2+w^2=1.
\end{equation}
となる。三次元球面は視覚化するのがより難しい。補助空間が今度は四次元になるからであるが、数学はそれほど難しくはない。

標準的な記法では\(n\)次元球面を\(\Omega_n\)と書く。したがって\(\Omega_1\)は一次元球面、\(\Omega_2\)は二次元球面、\(\Omega_3\)は三次元球面である。

\section{球面上の計量}

より有用な記述は内在的である。球面自身の上に住む座標によって球面を記述する。

単位一次元球面には一つの座標\(\phi\)で十分であり、\(0\leq \phi \leq 2\pi\)である。その計量は
\begin{equation}
ds^2=d\phi^2.
\end{equation}
である。これは
\begin{equation}
d\Omega_1^2=d\phi^2.
\end{equation}
とも書かれる。

単位二次元球面は一次元球面の列として記述される。一つの角度\(\phi\)は依然として各円の周りを回るが、第二の角度\(\theta\)は球面の一端から他端へと走る。その第二の角度は\(0\)から\(\pi\)まで変化する。\(\theta\)方向の運動は\(d\theta^2\)に寄与する。\(\phi\)方向の周りの運動は半径\(\sin\theta\)の円の上で起こる。したがって
\begin{equation}
d\Omega_2^2=d\theta^2+\sin^2\theta\,d\phi^2
\end{equation}
または
\begin{equation}
d\Omega_2^2=d\theta^2+\sin^2\theta\,d\Omega_1^2.
\end{equation}
となる。

三次元球面も同じようにして得られるが、今度は二次元球面の列としてである。もう一つ座標\(\alpha\)を導入する。すると
\begin{equation}
d\Omega_3^2=d\alpha^2+\sin^2\alpha\,d\Omega_2^2.
\end{equation}
となる。展開すると、これは
\begin{equation}
d\Omega_3^2=d\alpha^2+\sin^2\alpha\left(d\theta^2+\sin^2\theta\,d\phi^2\right).
\end{equation}
である。ここで\(\alpha\)は一方のゼロサイズの端から他方の端まで走るので、\(0\leq \alpha \leq \pi\)である。\editorialnote{00において：17：24音声は一時的に``ゼロから2パイ''のように聞こえるが、周囲の構成と明示的な\(\sin^2\alpha\)因子は、新しい座標が低次元球面の族の一方のゼロサイズの端から他方の端まで走ることを要求する。}

一般的な再帰的形式は
\begin{equation}
d\Omega_n^2=d\alpha^2+\sin^2\alpha\,d\Omega_{n-1}^2.
\end{equation}
である。

この観点を採用すれば、埋め込みの描像は忘れてよい。一次元球面は一つの座標で、二次元球面は二つで、三次元球面は三つで記述され、以下同様である。幾何学はそれ自体で成り立つ。

この幾何学が特別なのは、それが一様であるからである。\(\sin^2\theta\)を両端で依然として消える別の関数に置き換えて計量を変形することもできよう。そうすれば何らかの変形された閉曲面、たとえば楕円体が得られるであろう。しかしそれはもはや一様ではなくなる。球面が一様であるのは、どの点も他のどの点とも同じであるからである。

単位半径ではなく半径\(A\)の球面が欲しければ、単位計量に\(A^2\)を掛ける：
\begin{equation}
ds^2=A^2\,d\Omega_3^2.
\end{equation}
二乗が必要なのは、計量が距離の二乗を測るからである。

ちょっとした余談はゼロ次元球面に関するものである。一歩後ろに外挿すると
\begin{equation}
x^2=1,
\end{equation}
が得られ、その解は二点\(x=\pm 1\)である。したがって単位ゼロ次元球面は単に二点である。

\section{球状FRW宇宙}

今や関心は空間だけではなく時空にある。関連する描像は、半径が時間とともに変化する三次元球面である。

一つの警告が全体を通じて重要である：時間は異なる。膨張宇宙を、時間がもう一つの埋め込み座標の役割を果たす高次元球面として考えてはならない。

それを念頭に置くと、時空間隔は
\begin{equation}
d\tau^2=dt^2-a(t)^2\,d\Omega_3^2.
\end{equation}
である。これは空間的に球状な宇宙に対するFRW計量である。静的球面と比べて唯一の新しい要素は、半径が時間依存関数\(a(t)\)になったことである。

座標を、我々が一方の極にいるように選ぶと便利である。その点に特別な意味はなく、単なる約束事にすぎない。すると\(\alpha\)は、単位球面上で我々から離れる向きの動径座標と考えることができる。これはキロメートル単位の動径距離ではない。単位球面上の無次元の角度尺度であるが、幾何学においては動径距離の役割を果たす。

銀河や恒星のような局所的な塊を別にすれば、これが空間的に球対称なFRW宇宙の幾何学の全体である。その空間曲率は正であり、符号の約束は
\begin{equation}
k=+1.
\end{equation}
と書かれる。``何か''自体が1に足されるわけではない。これは単に空間曲率が正であることを意味する。

\section{球面から外を見る}

そのような球面を、その上の一点から外を見たとき、どのように見えるだろうか。ここでステレオ投影が登場する。

無限平面に南極で接する2次元球面を取る。球面上の一点を選ぶ。北極からその点を通る直線を引き、平面に当たるまで延ばす。このようにして、球面上の各点は平面上の一点に写される。南極は平面の原点に写る。北極は例外点である。すなわち、北極に異なる方向から近づくと、その像は方向に依存して無限遠へと走り去る。したがって、北極は無限遠点として扱われる。

等しい``大陸''を球面上の異なる位置に置くと、平面上の像は北極に近づくにつれて大きくなる。したがって、南極に置かれた観測者は、遠方の天体が平坦空間におけるよりも大きな角を張るのを見るであろう。正の曲率では、既知の天体はユークリッド幾何学に比べて大きすぎるように見える。

これが基本的な観測上の要点である。曲率は、距離、大きさ、角の大きさの関係から推測される。銀河自体は非常に深い距離に対する良い標準尺ではないが、論理はまさにそのとおりである。

\begin{classroomqa}
\lecturetimestamp{00:45:54}
\audiencequestion{角の大きさだけでなく、距離も知る必要がありますか。}
\lecturerresponse{はい。角度だけでは十分ではない。角度と距離の両方が必要である。}
\end{classroomqa}

\begin{classroomqa}
\lecturetimestamp{00:46:18}
\audiencequestion{球面上の円は円として投影されますか。}
\lecturerresponse{はい。円は円に投影される。これはステレオ投影の特別な特徴の一つである。}
\end{classroomqa}

\begin{classroomqa}
\lecturetimestamp{00:47:02}
\audiencequestion{大陸は投影の下で同じ形を保ちますか。}
\lecturerresponse{二つの異なる主張がある。一つは、円が円に投影されるということである。もう一つは、角度が保たれるため、十分に小さな図形はその局所的な形を保つという、より弱い主張である。この写像は等角写像である。すなわち、角度は保つが、相対的な大きさは保たない。}
\end{classroomqa}

\section{一様な負の曲率}

一様に負に曲がった場合は、その計量を通して導入される。

球面の負曲率版は双曲空間である。以前と同じ再帰的な精神で、その計量は\(\sin\alpha\)を\(\sinh\alpha\)に置き換えることによって得られる。
\begin{equation}
d\mathcal H_n^2=d\alpha^2+\sinh^2\alpha\,d\Omega_{n-1}^2.
\end{equation}
低次元の角度部分は単位球面計量\(d\Omega_{n-1}^2\)のままである。

\begin{classroomqa}
\lecturetimestamp{00:55:51}
\audiencequestion{計量において、その低次元の微分は依然として球面のものですか。}
\lecturerresponse{はい。その低次元の部分は通常の単位球面計量である。変化はその前の\(\sinh^2\alpha\)因子にある。}
\end{classroomqa}

双曲線正弦は
\begin{equation}
\sinh\alpha=\frac{e^\alpha-e^{-\alpha}}{2}.
\end{equation}
である。マイナス符号が重要である。講義では、学生によって直ちに訂正される。

座標の範囲にもう一つ決定的な違いがある。球面の場合、新しい座標は一方の端から他方の端まで進み、\(\pi\)で止まる。ここでは、
\begin{equation}
0\leq \alpha<\infty.
\end{equation}
である。通常の正弦は増加した後、再びゼロに戻る。双曲線正弦は\(\alpha=0\)の近くでは同じように始まるが、その後は指数関数的に増大する。

これは、我々の周りの低次元球面が、外側へ進むにつれてどんどん大きくなることを意味する。この幾何学では、半径が増大するにつれて空はまず大きくなり、その後縮むのではない。成長し続けるのである。

3次元空間では、対応するFRW形式は
\begin{equation}
d\tau^2=dt^2-a(t)^2\,d\mathcal H_3^2,
\end{equation}
であり、符号のラベルは
\begin{equation}
k=-1.
\end{equation}
である。

\section{円盤の絵とその意味}

負の曲率は、エッシャーによって有名になった円盤の絵、特に\emph{限界円第四番}において直観的になる。ステレオ投影との類推は有用であるが、振る舞いは逆である。

球面の場合、射影は無限に見えるが、実際の幾何学は有限である。負の曲率の場合、円盤の描像は有限に見えるが、実際の幾何学は無限である。等しい``大陸''や``銀河''は、境界に向かって進むにつれてますます小さく見えるが、幾何学的にはそれらはすべて同じ大きさである。距離を、いくつの等しい断片を横切るかを数えることによって測るならば、境界は無限に遠い。

したがって、負の曲率における既知の物体は、平坦空間に対して小さすぎるように見える。それは小さすぎる角度を張る。

幾何学は依然として一様である。描像は外側の断片を微小に見せるが、別の位置へ移動したどの観測者も同じ内在的幾何学を見るであろう。

\begin{classroomqa}
\lecturetimestamp{01:06:42}
\audiencequestion{それらの円のいずれかの中にいる観測者は、同様のパターンを見るであろうか。}
\lecturerresponse{はい、まったく同じものを見る。空間は一様である。}
\end{classroomqa}

\begin{classroomqa}
\lecturetimestamp{01:07:03}
\audiencequestion{球面の場合には遠くのものは大きすぎるように見えるが、ここでは小さすぎるように見えるのか。}
\lecturerresponse{はい。正の曲率では、既知の物体は平坦空間に対して大きすぎる角度を張る。負の曲率では、それは小さすぎる角度を張る。}
\end{classroomqa}

\begin{classroomqa}
\lecturetimestamp{01:08:09}
\audiencequestion{ステレオ投影的な射影の関連性とは、天文学者が観測からそれを再構成できるということか。}
\lecturerresponse{はい。天文学者は事実上、空の三角形を測量する測量士である。当面の幾何学的な点については、時間依存性を脇に置き、角度と距離に焦点を当てることができる。}
\end{classroomqa}

有用な局所的な描像もある。負の曲率の小さな一片は、通常の三次元空間における鞍面として描くことができる。しかし、そのように描くことができるのは小さな一片だけであり、双曲平面全体を三次元ユークリッド空間における素直な通常の曲面として大域的に埋め込むことはできない。

\section{平坦な場合}

正の曲率の場合と負の曲率の場合の間に、平坦空間が位置する。それを極限と考えることもできるが、単に通常の平坦なユークリッド空間と考えるだけで十分である。

そのFRW形式は
\begin{equation}
d\tau^2=dt^2-a(t)^2\left(dx^2+dy^2+dz^2\right),
\end{equation}
であり、
\begin{equation}
k=0.
\end{equation}
である。

これらが三つの一様なFRW幾何学である。
\begin{align}
k=+1 &: \text{positive spatial curvature},\\
k=0 &: \text{flat space},\\
k=-1 &: \text{negative spatial curvature}.
\end{align}

\begin{classroomqa}
\lecturetimestamp{01:15:07}
\audiencequestion{幾何学を平面に射影し、信号を追跡するならば、射影された光の速さは場所ごとに変化するであろうか。}
\lecturerresponse{はい、射影上では変化するであろう。球面射影では、外側へ行くほどますます速く進むように見えるであろう。負の曲率の円盤では、縁の近くでますます遅く進むように見えるであろう。その変化は描像に属するものであり、背後にある物理に属するものではない。}
\end{classroomqa}

\begin{classroomqa}
\lecturetimestamp{01:17:05}
\audiencequestion{実際に、このように遠くの物体を見ることを用いて平坦性を決定したのか。}
\lecturerresponse{はい、しかし実際に用いられるのは角度の情報である。射影上で光がどれほど速く動くように見えるかを測定しようとすることによって平坦性を決定するのではない。}
\end{classroomqa}

\section{宇宙論的方程式に戻る}

休憩の後には、すでに学んだ宇宙論的方程式とアインシュタイン方程式との関係が来る。この時点で、黒板にはまだ平坦な空間計量の断片が残っている。

\begin{figure}[htbp]
\centering
\includegraphics[width=0.72\textwidth]{lecture_05_figure_02.png}
\caption{平坦なFRW計量の断片。}
\lectureframe{01:18:18}
\end{figure}

まず非常に短いニュートン的な復習を行う。我々を中心に置く。ハッブル則に従って遠ざかる銀河を見る。球対称な分布に対するニュートンの定理により、銀河の距離に一致する半径の球の内側の質量だけが、それに働く力に寄与する。すると三つの可能性がある。銀河が脱出速度を超えて運動するか、脱出速度を下回るか、ちょうど脱出速度にあるかである。

次に想起されるのは、通常の非相対論的物質では密度が質量を\(a^3\)で割ったもののように振る舞い、放射では\(1/a^4\)のように振る舞い、真空エネルギーでは球が変化しても変化しないということである。

以前のニュートン的な取り扱いでは、宇宙論的方程式には密度の項と、脱出速度の場合にちょうど位置しなければ\(1/a^2\)に比例する追加の定数項があった。以前の議論では、その項を主にゼロに置いていた。経験的にはそれは小さいが、概念的にはそれを省く理由はない。

\section{アインシュタイン方程式と\(00\)成分}

次の段階はアインシュタインの場の方程式である。左辺には時空の曲率から組み立てられた幾何学が立つ。右辺には重力場の源であるエネルギー・運動量テンソルが立つ。講義の規格化では次のように書く。
\begin{equation}
R^{\mu\nu}-\frac{1}{2}g^{\mu\nu}R=4\pi G\,T^{\mu\nu}.
\end{equation}
\editorialnote{可視の板書行では因子\(\tfrac12\)は判読できないが、01:33:36のあたりで口頭の訂正がそれを復元する。ここでの方程式はその訂正を保ちつつ、講義の\(4\pi G\)規格化を保っている。}

ここでは十個の独立成分すべてが必要なわけではない。この段階までにすでにあまりに多くが仮定されている---一様性、FRW形式、そして時間のみに依存する単一のスケール因子---ので、残る任意関数はただ一つ、すなわち\(a(t)\)だけである。選ぶべき自然な成分は時間-時間成分である。なぜなら\(T^{00}\)はエネルギー密度\(\rho\)だからである。

\begin{figure}[htbp]
\centering
\includegraphics[width=0.62\textwidth]{lecture_05_figure_05.png}
\caption{アインシュタインの場の方程式とその\(00\)成分。}
\lectureframe{01:30:54}
\end{figure}

したがって焦点は
\begin{equation}
R^{00}-\frac{1}{2}g^{00}R=4\pi G\,\rho.
\end{equation}
である。

右辺はエネルギー密度である。左辺は幾何学である。

\begin{classroomqa}
\lecturetimestamp{01:39:09}
\audiencequestion{エネルギー・運動量テンソルの他の成分はどうなるのか。}
\lecturerresponse{それらは一般にはゼロにならない。空間-空間成分は圧力を表しうる。物質優勢の場合には圧力はゼロであり、放射優勢の場合にはエネルギー密度の三分の一であり、真空エネルギーの場合には\(P=w\rho\)と書く。ただし\(w=-1\)である。}
\end{classroomqa}

\begin{classroomqa}
\lecturetimestamp{01:40:24}
\audiencequestion{では、物質しかなければ\(T^{\mu\nu}\)の他の成分はゼロになるのか。}
\lecturerresponse{圧力のない物質については、そうだ、圧力成分はゼロになる。}
\end{classroomqa}

\section{アインシュタインの計算がもたらすもの}

完全な曲率の計算は板書では行われない。重要なのは、FRW計量をアインシュタインテンソルに代入したときに現れるものである。

一方の種類の項は、スケール因子が時間に依存するという事実だけから生じる。空間が平坦であっても、時間に依存するスケール因子は依然として時空の曲率を与える。その寄与は次に比例する。
\begin{equation}
\left(\frac{\dot a}{a}\right)^2.
\end{equation}

第二の種類の項は空間曲率そのものから生じる。スケール因子が時間的に凍結されていても、曲がった空間は依然として時空の曲率に寄与する。単純な一様に曲がった幾何学では、曲率は曲率半径の二乗の逆数に比例するので、次に比例する項に寄与する。
\begin{equation}
\frac{k}{a^2},
\end{equation}
球面の場合には正の符号、双曲的な場合には負の符号、平坦な空間の場合にはゼロである。

\begin{classroomqa}
\lecturetimestamp{01:36:53}
\audiencequestion{曲率半径のような因子もまたあるべきではないか。}
\lecturerresponse{それはすでに\(1/a^2\)の因子である。これらの単純な一様に曲がった幾何学では、曲率はちょうど曲率半径の二乗の逆数である。}
\end{classroomqa}

係数の帳尻合わせはその場で行われ、部分的には以前のニュートン的な答えと照合される。終着点は\footnote{である。この表示されたフリードマン形式は、口頭で自己訂正された係数の議論と、以前のニュートン的結果との比較を01:32:35--01:36:53のあたりで一つにまとめている。}
\begin{equation}
\left(\frac{\dot a}{a}\right)^2+\frac{k}{a^2}=\frac{8\pi G}{3}\rho,
\end{equation}
あるいは等価に
\begin{equation}
\left(\frac{\dot a}{a}\right)^2=\frac{8\pi G}{3}\rho-\frac{k}{a^2}.
\end{equation}

これは三つのFRW幾何学を以前の宇宙論的力学と結びつける形式である。

\begin{classroomqa}
\lecturetimestamp{01:37:52}
\audiencequestion{\(a\)を再定義すれば\(k\)は任意の値を取りうるのか。}
\lecturerresponse{平坦な場合には\(a\)に対する一意な自然なスケールが存在しないので、それを再スケールできるが、\(k\)はゼロのままである。球面の場合と負に曲がった場合には自然な曲率スケールが存在する。\(a(t)\)をその自然な曲率半径と取れば、\(k\)はちょうど\(+1\)または\(-1\)である。}
\end{classroomqa}

\begin{classroomqa}
\lecturetimestamp{01:41:00}
\audiencequestion{その最終方程式で\(k\)は\(+1\)、\(-1\)、\(0\)のいずれに等しいのか。}
\lecturerresponse{そのとおり。球面の場合には\(+1\)、双曲的な場合には\(-1\)、平坦な空間の場合には\(0\)である。}
\end{classroomqa}

ニュートンの扱いは我々を誤った道に導かなかった。それは同じ基本方程式を与えた。一般相対性理論が付け加えるのは解釈である。ニュートン的な観点からは、余分な定数は脱出速度を上回る場合、脱出速度を下回る場合、そしてその境界の場合を区別する。アインシュタイン的な観点からは、同じ符号が空間そのものの幾何学を教えてくれる。

\begin{classroomqa}
\lecturetimestamp{01:42:46}
\audiencequestion{宇宙定数を導入する前には、これはどのように成り立っていたのか。}
\lecturerresponse{宇宙定数の前には、脱出速度の振る舞いと空間の幾何学との間に密接な関係があった。\(k=+1\)ならば、空間は正の曲率をもち、運動は脱出速度を下回るので、宇宙は膨張した後に再収縮しうる。\(k=-1\)ならば、空間は負の曲率をもち、運動は脱出速度を上回るので、宇宙は膨張し続ける。\(k=0\)ならば、境界的なぎりぎりの場合となる。}
\end{classroomqa}

したがって、同じ符号\(k\)は二通りに読める。すなわち、運動の動力学的な分類として、また空間の曲率の符号としてである。
